\section{マクロ経済スライド}

\begin{fillblank}[年金1 2021 問題1(5)\textcircled{2}]
    \noindent \textbf{◯国民年金法 附則(平成一六年六月一一日法律第一〇四号)}
  \begin{lawstructure}
    \begin{lawarticle}{第二条}{給付水準の下限}
      国民年金法による年金たる給付及び厚生年金保険法による年金たる保険給付については、
      第一号に掲げる額と第二号に掲げる額とを合算して得た額の第三号に掲げる額に対する比率が
      （　　A　　）を上回ることとなるような給付水準を将来にわたり確保するものとする。

      \lawitem{一}
      当該年度における国民年金法による老齢基礎年金の額（当該年度において（　　B　　）に達し、
      かつ、保険料納付済期間の月数が（　　C　　）である受給権者について計算される額とする。）を
      当該年度の前年度までの標準報酬平均額（厚生年金保険法第四十三条の二第一項第二号イに
      規定する標準報酬平均額をいう。）の推移を勘案して調整した額を（　　D　　）額に（　　E　　）
      得た額に相当する額

      \lawitem{二}
      当該年度における厚生年金保険法による老齢厚生年金の額
      （当該年度の前年度における男子である同法による被保険者（次号において「男子被保険者」という。）
      の平均的な標準報酬額（同法による標準報酬月額と標準賞与額の総額を（　　D　　）得た額とを
      合算して得た額をいう。次号において同じ。）に相当する額に当該年度の前年度に属する月の
      標準報酬月額又は標準賞与額に係る再評価率（同法第四十三条第一項に規定する再評価率をいい、
      当該年度に（　　B　　）に達する受給権者に適用されるものとする。）を乗じて得た額を平均標準報酬額
      とし、被保険者期間の月数を（　　C　　）として同項の規定の例により計算した額とする。）を
      （　　D　　）得た額に相当する額

      \lawitem{三}      
      当該年度の前年度における男子被保険者の平均的な標準報酬額に相当する額から当該額に係る公租公課の額を控除して得た額に相当する額
    
      \lawparagraph{2〜3　（略）}
    \end{lawarticle}
  \end{lawstructure}
\end{fillblank}


\begin{answerbox}
  \begin{enumerate}[A: ]
    \item 百分の五十
    \item 六十五歳
    \item 四百八十
    \item 十二で除して
    \item 二を乗じて
  \end{enumerate}
\end{answerbox}

\begin{fillblank}[]
  2021 年（令和3 年）4 月1 日から、賃金スライドおよび物価スライドのルールが変更され
る。変更後は、例えば、名目手取り賃金変動率が1.0％減の0.990、物価変動率が 1.0％増の
1.010 のときの老齢基礎年金の改定率は、新規裁定年金は（　　A　　）を基準として、既裁定年
金は（　　B　　）を基準として、それぞれ改定される。
\end{fillblank}

\begin{answerbox}
  \begin{enumerate}[A: ]
    \item 0.990
    \item 0.990
  \end{enumerate}
\end{answerbox}

\begin{explanationbox}
  原則：新規裁定年金は賃金変動率、既裁定年金は物価変動率をベースに改訂

  物価変動率＞賃金変動率の場合：既裁定年金も賃金変動率で改訂
\end{explanationbox}

\begin{fillblank}[年金1 H28 問題1(2)H]
   平成 12 年度から平成 14 年度にかけて実施された物価下落に対応した年金額の特例措置によ
る給付水準（特例水準）と、本来の給付水準との乖離分は、平成 23 年度には最大（　　A　　）％まで
拡大した。
\end{fillblank}


\begin{answerbox}
  \begin{enumerate}[A: ]
    \item 2.5
  \end{enumerate}
\end{answerbox}

%毎年変える
\begin{fillblank}[自作 2026/11/3]
  令和8年度の年金額改定に使用する参考指標及び年金額の改定率について以下の空欄を埋めなさい。

  \begin{itemize}
    \item （令和7年の）物価変動率：（　　A　　）％
    \item （令和7年の）名目手取り賃金変動率：（　　B　　） ％ 
    \begin{flalign*}
      &（令和7年の）名目手取り賃金変動率 \\
      &= \text{実質賃金変動率 (令和4〜6年度の平均)}　（　　\text{C}　　） ％ \\
      &+ \text{物価変動率 (令和7年)}　（　　\text{A}　　） ％ \\
      &+ \text{可処分所得割合変化率 (令和5年度)}　（　　\text{D}　　） ％ 
    \end{flalign*}

    \item （令和7年の）マクロ経済スライドによるスライド調整率：（　　E　　） ％ 
    \begin{flalign*}
      &（令和7年の）マクロ経済スライドによるスライド調整率 \\
      &= \text{公的年金被保険者総数の変動率（令和4〜6年度の平均）}　（　　\text{F}　　） ％ \\
      &+ \text{平均余命の伸び率（定率）}　（　　\text{G}　　） ％
    \end{flalign*}

    \item 前年度以前までのマクロ経済スライドの未調整分：（　　H　　）％
    \item （新規裁定者・既裁定者の）令和8年度の国民年金(基礎年金)年金額の改定率：（　　I　　）％
    \item （新規裁定者・既裁定者の）令和8年度の厚生年金(報酬比例部分)年金額の改定率：（　　J　　）％
  \end{itemize}
\end{fillblank}


\begin{answerbox}
  \begin{enumerate}[A: ]
    \item 3.2
    \item 2.1
    \item -1.1
    \item 0.0
    \item -0.2
    \item 0.1
    \item -0.3
    \item 0.0
    \item $2.1-0.2 = 1.9$（物価＞賃金なので、既裁定者も賃金変動率で改定）
    \item 厚生年金の調整率は$1/3$に緩やかにするので、スライド調整率は$-0.2\times1/3=-0.06666\cdots=-0.1$。\\
      よって、$2.1-0.1=2.0$。
  \end{enumerate}
\end{answerbox}

\begin{explanationbox}
  令和８年度の年金額改定について

  \url{https://www.mhlw.go.jp/stf/houdou/0000191631_00020.html}
\end{explanationbox}

\begin{shortanswer}[年金1 2018 問題2(5)\textcircled{1}]
  公的年金の財政検証においては、次の財政検証までに所得代替率が 50％を下回ると見込まれ
る場合には、給付水準調整の終了その他の措置を講ずるとともに、給付および負担の在り方に
ついて検討を行い、所要の措置を講ずるものとされている。ここで用いられる所得代替率の算
定方法を簡記しなさい。
\end{shortanswer}

\begin{answerbox}
  「厚生年金の標準的な受給世帯の新規裁定時の年金額」の「現役世代の男子の平均手取り収入額」に対する比率が所得代替率である。
\end{answerbox}

\begin{shortanswer}[年金1 H29 問題2(5)]
  「公的年金制度の持続可能性の向上を図るための国民年金法等の一部を改正する法律」（平成 28
年法律第 114 号）により年金額の改定ルールの見直しが実施されることに関し、次の(1)、(2)の設
問に答えなさい。
\begin{enumerate}[(1)]
  \item 平成 33 年 4 月施行の年金額改定ルールの見直しについて、その内容を簡記しなさい。
\item 平成 33 年 4 月施行の年金額改定ルールの見直しの理由として、平成 16 年以降の過去の物価・
賃金等の動向により、 平成 16 年改正時の想定と比べて将来見通しにどのようなズレが生じてき
たことが問題視されたのか、簡記しなさい。
\end{enumerate}

\end{shortanswer}

\begin{answerbox}
  \begin{enumerate}[(1)]
    \item 賃金変動が物価変動を下回る場合であって、かつ、賃金変動がマイナスの場合（物価＞０＞賃金、又は、０＞物価＞賃金）について、既裁定者及び新規裁定者の年金額を賃金変動率に応じて改定するように変更
    \item デフレ下での実質賃金下落により、現在の受給者世代の所得代替率が上昇する一方で、基礎年金部分で将来の受給者の所得代替率が低下する見通しとなったこと。
  \end{enumerate}
\end{answerbox}

\begin{explanationbox}
  昔の時事問題
\end{explanationbox}

\begin{shortanswer}[年金1 H28 問題2(4) 改題]
  公的年金に関する次の(1)～(2)の設問に答えなさい。
  \begin{enumerate}[(1)]
    \item 平成 6 年改正により標準報酬の再評価の方法が変更されているが、その変更内容を簡記しな
さい。
    \item 平成 16 年改正により導入されたマクロ経済スライドにおけるスライド調整率とは、何を反映
するためのものか、簡記しなさい。
  \end{enumerate}
\end{shortanswer}

\begin{answerbox}
  \begin{enumerate}[(1)]
    \item 「一人当たりグロスの賃金上昇率による方法」から「一人当たり可処分所得時の上昇率による方法」へと変更になった
    \item 「公的年金全体の被保険者の減少」と「平均的な年金受給期間の伸び」を年金改定率に反映
  \end{enumerate}  
\end{answerbox}


\begin{shortanswer}[年金1 2023 問題2(3)ア]
  国民年金及び厚生年金の年金額改定に関して、次の（ａ）、（ｂ）それぞれの算出方法を簡潔
に入力しなさい。ただし、具体的な数値には触れなくてよい。（（ａ）、（ｂ）それぞれ150 字
以内）

\begin{enumerate}[(a)]
  \item 名目手取り賃金変動率
  \item マクロ経済スライドによるスライド調整率（前年度までの未調整分については触れなくて
よい）
\end{enumerate}

\end{shortanswer}


\begin{answerbox}
\begin{enumerate}[(a)]
  \item ２年度前から４年度前までの３年度平均の実質賃金変動率に前年の物価変動率と３年度前
の可処分所得割合変化率を乗じて算出
  \item ２年度前から４年度前までの３年度平均の公的年金全体の被保険者数の減少率に平均余命
の伸びを勘案した一定率を加えて算出
\end{enumerate}
 
\end{answerbox}

\begin{shortanswer}[年金1 2024 問題2(3)ア]
  平成 16 年年金制度改正で枠組みが定められた公的年金の給付と財源の関係について、簡潔に
入力しなさい。なお、「保険料」「国庫負担」「マクロ経済スライド」「所得代替率」をすべて用
いること。
\end{shortanswer}

\begin{answerbox}
  財源としては、保険料水準は固定され、基礎年金の国庫負担割合が1/2、また積立金が活用され
る。財源が固定される中で、少子高齢化が進行しても給付を賄えるよう、長期的に財政収支を均
衡させていくために、マクロ経済スライドによる年金の実質価値の引き下げが行われる。ただ
し、次回の財政検証までに所得代替率が 50\%を下回ると見込まれる場合には、調整期間の終了
について検討を行い、その結果に基づいて調整期間の終了その他の措置を講ずるものとされて
いる
\end{answerbox}


\begin{fillblank}[年金2 2021 問題1(4)]
  次は公的年金におけるマクロ経済スライドに関する記述である。以下の空欄を埋めなさい。
\\

  所得代替率は、マクロ経済スライドによる給付水準調整が行われなければ原則として
（　　A　　）。一方、マクロ経済スライドが発動されると、（　　B　　）。

しかしながら、デフレ経済が長引いたことにより特例水準が解消しなかったため、2004 年
改正で導入されたマクロ経済スライドが発動し給付水準の調整が行われたのは（　　C　　）年度の
0.9％の調整が初めてである。

なお、2004 年導入時からマクロ経済スライドが発動されなかった間に所得代替率は若干
（　　D　　）した。この所得代替率の（　　D　　）は、主に（　　E　　）で起こっている。

2016（平成28）年に公的年金制度の持続可能性の向上を図るための国民年金保険法等の一
部の法律の改正により、年金額の改訂ルールの見直しが行われた。マクロ経済スライドについ
ては、年金の（　　F　　）が前年度を下回らない措置を継続しつつ、賃金・物価上昇の範囲内で前
年度までの未調整分を含めて調整する仕組みが導入され、マクロ経済スライドの強化が図られ
た。

あわせて賃金変動が物価変動を下回る場合には（　　G　　）変動にあわせて年金額を改定するこ
ととし、年金財政の支え手である現役世代の負担能力に応じた年金給付とする仕組みが導入さ
れた。
\end{fillblank}

\begin{answerbox}
  \begin{enumerate}[A: ]
    \item 変わらない
    \item 低下する
    \item 2015
    \item 上昇
    \item 基礎年金
    \item 名目額
    \item 賃金
  \end{enumerate}
\end{answerbox}


\begin{fillblank}[年金2 2023 問題1(4)]
  2016(平成28)年成立の「公的年金制度の持続可能性の向上を図るための国民年金法等の一
部を改正する法律」により年金額改定ルールが改正された。厚生労働省「2019（令和元）年財
政検証レポート」を参考にして、この法律に基づく年金額改定ルールのうちマクロ経済スライ
ドによる調整前の取扱いを整理した下表について以下の空欄を埋めなさい。

  \begin{table}[htbp]
    \centering
    \begin{tabular}{|c|c|c|} \hline 
      \multirow{2}{*}{賃金と物価の関係} &   \multicolumn{2}{|c|}{年金額改定（スライド）のルール} \\ 
      \cline{2-3} & 新規裁定者（（　　A　　）歳以下） & 既裁定者（（　　※　　）歳以上） \\ \hline
      $賃金 > 物価$ & （　　※　　） & （　　B　　） \\ \hline
      $物価 > 賃金 > 0$ & （　　※　　） & （　　C　　） \\ \hline
      $物価 > 0 > 賃金$ & （　　D　　） & （　　※　　） \\ \hline
      $0 > 物価 > 賃金$& （　　※　　） & （　　E　　） \\ \hline
    \end{tabular}
  \end{table}

  （※）は出題の都合上省略している。

  （補足）賃金：名目手取り賃金変動率、物価：物価変動率
\end{fillblank}

\begin{answerbox}
  \begin{enumerate}[A: ]
    \item 67
    \item 物価変動率を基準に改定
    \item 名目手取り賃金変動率を基準に改定
    \item 名目手取り賃金変動率を基準に改定
    \item 名目手取り賃金変動率を基準に改定
  \end{enumerate}
\end{answerbox}

\begin{shortanswer}[数理人試験 2025 問題3 設問2]
  第21回社会保障審議会年金部会資料1「基礎年金のマクロ経済スライドによる給付調整の早期終
了（マクロ経済スライドの調整期間の一致）について」において記載がある、マクロ経済スライド
の調整期間の一致により将来の給付水準が上昇する要因を2つ簡記せよ。
\end{shortanswer}

\begin{answerbox}
  \begin{itemize}
    \item 基礎（１階）の水準上昇に伴う国庫負担の増加で給付が純増するため。
\item 比例（２階）の給付調整が進むことで足下の受給世代の比例（２階）の財源が将来の
受給世代の基礎（１階）の給付に充てられ、世代間の財源移転が行われるため。
  \end{itemize}
\end{answerbox}

\newpage